\documentclass{article}
\usepackage{amsmath,amssymb}
\begin{document}

\section*{EF–CFI Transfer Formalization}

\textbf{State Space:}
\[
S_t = (\bar a_t,\bar b_t),\quad
\bar a_t\in V(G_0)^k,\;\bar b_t\in V(G_1)^k,
\]
with partial isomorphism invariant:
\[
\forall i,j:\;
(a_i=a_j)\Leftrightarrow(b_i=b_j),\quad
E(a_i,a_j)\Leftrightarrow E(b_i,b_j).
\]

\textbf{Transition Function:}
\[
T:\;S_t\times (V(G_0)\cup V(G_1))\to S_{t+1}
\]
defined by extending tuples and truncating to last \(k\) entries.

\textbf{Extension Lemma (Tree Neighborhoods):}
\[
\mathrm{girth}(B)>2r\Rightarrow
B_r(v)\cong T_r(v)
\Rightarrow
\exists!\;\text{isomorphism extension preserving adjacency}.
\]

\textbf{Duplicator Response:}
\[
\forall \text{ Spoiler move }x\in G_0,\;
\exists y\in G_1:
\]
\[
B_r(x)\cong B_r(y)
\;\land\;
S_{t+1}\text{ remains partial isomorphism}.
\]

\textbf{CFI Symmetry Invariance:}
\[
\text{local automorphisms act transitively on gadget states}
\Rightarrow
\text{parity indistinguishable locally}.
\]

\textbf{Induction:}
\[
(\forall t\le k)\;S_t\text{ is partial isomorphism}
\Rightarrow
\text{Duplicator wins }k\text{-round EF game}.
\]

\textbf{Conclusion:}
\[
G_0\equiv_{FO^k}G_1.
\]

\textbf{Status: OPEN (requires formal proof of uniqueness and automorphism action).}

\end{document}
